\subsection{Đặt bài toán}

\subsubsection{Mục tiêu}

Trong các tổ chức quy mô vừa và nhỏ như trường học, doanh nghiệp hay cộng đồng địa phương, nhu cầu tổ chức các cuộc bỏ phiếu nội bộ ngày càng phổ biến: bầu ban chấp hành, lấy ý kiến về một chủ trương quan trọng, bình chọn cá nhân tiêu biểu, hay đơn giản là quyết định một lựa chọn chung của tập thể. Các cuộc bầu cử dạng này có chung một đặc điểm là phạm vi tham gia khép kín -- chỉ những người thuộc tổ chức và đã được xác định trước mới có quyền bỏ phiếu -- đồng thời thường được điều hành bởi một bộ phận cụ thể chịu trách nhiệm tổ chức và công bố kết quả.

Hệ thống bỏ phiếu được trình bày trong đồ án này hướng đến đúng đối tượng nêu trên: các cuộc bầu cử nội bộ trong một tổ chức nhỏ, trong đó có một quản trị viên đại diện tổ chức cấp tài khoản định danh cho từng cử tri trước khi cuộc bầu cử bắt đầu. Nhờ vậy, hệ thống biết rõ tập cử tri hợp lệ là ai và không phải xử lý vấn đề định danh ở quy mô toàn xã hội -- vốn là một bài toán rất phức tạp khi triển khai bầu cử công khai trên Internet.

Mục tiêu tổng thể của hệ thống là tạo ra một môi trường bỏ phiếu điện tử thuận tiện cho người dùng, đồng thời vẫn bảo vệ được hai giá trị cốt lõi mà bất kỳ cuộc bầu cử nào cũng cần: \textbf{tính bí mật} của lá phiếu và \textbf{tính minh bạch} của kết quả. Người dùng có thể tham gia bỏ phiếu từ thiết bị cá nhân mà không phải đến tận nơi tổ chức như mô hình truyền thống, trong khi tổ chức vẫn có khả năng kiểm soát danh sách cử tri, thời điểm mở và đóng cuộc bầu cử, cũng như kiểm chứng tính trung thực của kết quả sau khi kết thúc.

Cụ thể, hệ thống được thiết kế nhằm đạt được \textbf{ba mục tiêu cốt lõi} sau:

\begin{itemize}
    \item \textbf{Bỏ phiếu duy nhất một lần:} Mỗi cử tri hợp lệ chỉ có thể bỏ phiếu đúng một lần trong một cuộc bầu cử. Hệ thống ngăn chặn các hành động bỏ phiếu nhiều lần dù do cố ý hay do lỗi thao tác.

    \item \textbf{Bí mật lá phiếu:} Lựa chọn của từng cử tri được giữ bí mật tuyệt đối. Không một thành phần nào trong hệ thống, kể cả quản trị viên hay máy chủ vận hành có thể biết được một cử tri cụ thể đã chọn ứng viên nào.

    \item \textbf{Tính kiểm chứng của kết quả:} Kết quả bầu cử có thể được xác minh một cách độc lập, để các cử tri và các bên quan tâm tự xác minh rằng phiếu bầu đã được tính và tổng số phiếu công bố là chính xác.
\end{itemize}

\subsubsection{Yêu cầu chức năng}

Từ mục tiêu đã đặt ra ở phần trên, hệ thống được phân tách thành ba nhóm chức năng chính tương ứng với ba vai trò người dùng: \textbf{quản trị viên} đại diện cho tổ chức, \textbf{cử tri} tham gia bỏ phiếu và \textbf{khách} là cử tri tương tác ở trạng thái ẩn danh hoặc bất kì bên thứ ba khác muốn kiểm chứng. Bên cạnh đó, hệ thống còn cần một số chức năng nền tảng nhằm bảo đảm các tính chất bảo mật cốt lõi mà mục tiêu đã đề ra.

\begin{itemize}
    \item \textbf{Chức năng dành cho quản trị viên:} Quản trị viên là người chịu trách nhiệm tổ chức cuộc bầu cử. Trước khi cuộc bầu cử bắt đầu, quản trị viên có khả năng tạo cuộc bầu cử mới, quản lý danh sách cử tri hợp lệ cùng với việc cấp tài khoản định danh tương ứng cho từng người. Trong quá trình diễn ra, quản trị viên có quyền mở phiên bỏ phiếu để các cử tri có thể tham gia và đóng phiên khi đến thời điểm kết thúc. Sau khi đóng phiên, quá trình kiểm phiếu và công bố kết quả chính thức được tự động kích hoạt.

    \item \textbf{Chức năng dành cho cử tri:} Mỗi cử tri sử dụng tài khoản đã được cấp để đăng nhập vào hệ thống. Sau khi đăng nhập, cử tri có thể xem được thông tin của cuộc bầu cử, danh sách ứng viên, lựa chọn ứng viên mà mình muốn bầu và gửi lá phiếu đi. Khi quá trình gửi phiếu hoàn tất, cử tri nhận được một biên lai làm bằng chứng cá nhân rằng phiếu của mình đã được hệ thống tiếp nhận và lưu giữ biên lai để sử dụng trong các chức năng kiểm chứng công khai.

    \item \textbf{Chức năng dành cho khách:} Khách là cử tri ở trạng thái ẩn danh hoặc bất kỳ bên thứ ba nào muốn kiểm chứng tính minh bạch của cuộc bầu cử. Sau khi cuộc bầu cử kết thúc, khách có thể tiết lộ lá phiếu ẩn danh và sử dụng biên lai cá nhân để xác minh phiếu đã được tính vào kết quả, đối chiếu dữ liệu giữa cơ sở dữ liệu và blockchain, đồng thời kiểm toán kết quả và trạng thái cam kết gốc Merkle. Các thao tác này không yêu cầu phiên xác thực mà chỉ sử dụng thông tin phiếu cần thiết, qua đó bảo đảm sự tách biệt giữa danh tính cử tri và nội dung lá phiếu, đồng thời cho phép mọi bên kiểm tra độc lập tính toàn vẹn của cuộc bầu cử.
    
    \item \textbf{Xác nhận tính hợp lệ của phiếu bằng chữ ký mù tập thể:} Để vừa có thể xác nhận một lá phiếu là hợp lệ vừa không nhìn thấy nội dung lá phiếu, hệ thống áp dụng cơ chế \textbf{chữ ký mù tập thể}. Cụ thể, các bên có thẩm quyền sẽ ký xác nhận lên một dạng \textit{phiếu đã bị che} mà không biết nội dung thật bên trong. Nhờ đó, máy chủ và các thành phần trung gian không có đủ thông tin để liên kết một cử tri cụ thể với một lựa chọn ứng viên cụ thể, giảm rủi ro gian lận từ chính bên vận hành hệ thống.

    \item \textbf{Ghi nhận dữ liệu trên blockchain nội bộ:} Toàn bộ các phiếu hợp lệ và các sự kiện quan trọng của cuộc bầu cử (mở phiên, đóng phiên, công bố kết quả) đều được ghi nhận lên một \textbf{blockchain nội bộ} do tổ chức vận hành. Đặc tính bất biến của blockchain bảo đảm dữ liệu sau khi ghi gần như bất kha thi trong việc sửa đổi dữ liệu, qua đó nâng cao tính toàn vẹn và độ tin cậy của kết quả cuối cùng.

\end{itemize}

Tóm lại, sự gắn kết chặt chẽ giữa các mục tiêu cốt lõi và yêu cầu chức năng đã tạo nên một khung giải pháp toàn diện cho bài toán bầu cử nội bộ. Từ bối cảnh thực tế của các tổ chức vừa và nhỏ, hệ thống đã cụ thể hóa ba nguyên tắc cốt lõi -- tính duy nhất của lá phiếu, sự bảo mật tuyệt đối và khả năng đối chiếu minh bạch -- thành các luồng nghiệp vụ thực tế. Sự kết hợp giữa các kỹ thuật mật mã như chữ ký mù tập thể \textit{(nhằm tách bạch danh tính khỏi nội dung phiếu)} và blockchain nội bộ \textit{(nhằm lưu vết bất biến)} đóng vai trò là xương sống công nghệ. Nhờ vậy, hệ thống không chỉ số hóa quy trình bầu cử một cách tiện lợi mà còn giảm rủi ro nghịch lý giữa việc bảo vệ quyền riêng tư cá nhân và đảm bảo tính minh bạch.

\subsubsection{Yêu cầu hệ thống}

Một hệ thống bầu cử điện tử đáng tin cậy phải đồng thời thỏa mãn năm tính chất sau đây. Các tính chất này không độc lập mà bổ trợ cho nhau, nếu thiếu một trong các tính chất này, mức độ tin cậy của kết quả bầu cử sẽ bị suy giảm đáng kể.

\begin{itemize}
    \item \textbf{Tính bí mật:} Kể cả quản trị viên hay bên vận hành hệ thống, không ai có thể biết một cử tri cụ thể đã bầu cho ứng viên nào. Bí mật của lá phiếu là nền tảng để loại trừ mọi tác động bên ngoài lên quyết định bỏ phiếu.

    \item \textbf{Tính duy nhất:} Mỗi cử tri hợp lệ chỉ được phép bỏ đúng một phiếu trong suốt một cuộc bầu cử. Hệ thống phải ngăn chặn mọi hình thức bỏ phiếu lặp lại, dù là cố tình hay do sự cố kỹ thuật.

    \item \textbf{Tính xác thực:} Chỉ những cử tri đã được tổ chức xác nhận trước mới có thể tham gia bỏ phiếu. Không một người ngoài hay tài khoản giả mạo nào có thể chen vào quá trình.

    \item \textbf{Tính không thể chối bỏ:} Một khi phiếu bầu đã được nộp và ghi nhận, không bên nào có thể phủ nhận sự kiện đó. Cả cử tri lẫn hệ thống đều không thể xóa bỏ hay làm giả dấu vết của một lá phiếu đã tồn tại.

    \item \textbf{Tính ẩn danh:} Ngay cả khi thu thập toàn bộ dữ liệu công khai của hệ thống, không thể truy ngược từ một lá phiếu về danh tính của người đã bỏ phiếu đó. Tính ẩn danh đi xa hơn tính bí mật: không chỉ giữ bí mật nội dung, mà còn hạn chế mọi mối liên hệ có thể phát hiện được giữa cử tri và lựa chọn của họ.
\end{itemize}

Để đảm bảo các tính chất trên trong thực tế, hệ thống áp dụng nguyên tắc \textbf{bảo vệ đa tầng} (\textit{defense in depth}). Thay vì phụ thuộc vào một cơ chế bảo vệ duy nhất, mỗi tầng trong kiến trúc đều có lớp kiểm soát riêng. Nhờ đó, nếu một tầng bị vượt qua, các tầng còn lại vẫn giữ vững toàn vẹn của cuộc bầu cử.

\begin{itemize}
    \item \textbf{Tầng giao tiếp mạng (HTTP):} Mọi yêu cầu từ trình duyệt của người dùng đều phải đi qua một cổng vào duy nhất. Tại đây, hệ thống tự động chặn các nguồn không được phép, giới hạn tần suất gửi yêu cầu để ngăn tấn công dồn dập, và bảo vệ các tiêu đề phản hồi để hạn chế lộ thông tin về cấu trúc bên trong.

    \item \textbf{Tầng xác thực và phân quyền:} Sau khi đăng nhập thành công, mỗi người dùng nhận một mã xác thực có thời hạn. Mọi thao tác tiếp theo đều phải trình mã này. Hệ thống kiểm tra không chỉ mã có hợp lệ hay không, mà còn kiểm tra người dùng có đúng vai trò được phép thực hiện thao tác đó hay không -- quản trị viên và cử tri có quyền hạn hoàn toàn khác nhau.

    \item \textbf{Tầng mật mã (chữ ký mù tập thể):} Lá phiếu được xử lý bằng một kỹ thuật mật mã đặc biệt: cử tri tự che nội dung lá phiếu trước khi gửi lên để ký xác nhận. Các bên ký chỉ xác nhận phiếu hợp lệ mà không nhìn thấy nội dung bên trong. Khóa ký được chia đều cho nhiều thành phần độc lập -- chỉ khi tất cả cùng tham gia mới tạo ra chữ ký hợp lệ -- loại trừ nguy cơ một bên duy nhất kiểm soát toàn bộ quá trình ký.

    \item \textbf{Tầng cơ sở dữ liệu:} Cơ sở dữ liệu áp đặt các ràng buộc: mỗi cử tri chỉ có một phiếu, mỗi phiếu là duy nhất, và mỗi lá phiếu chỉ được tiết lộ một lần. Bất kỳ hành động nào vi phạm các ràng buộc này đều bị từ chối tự động.

    \item \textbf{Tầng blockchain (Hyperledger Fabric):} Toàn bộ sự kiện quan trọng -- từng phiếu bầu, tổng hợp kết quả, và các mốc thời gian của cuộc bầu cử -- được ghi lên một sổ cái phân tán không thể sửa đổi. Bất kỳ ai cũng có thể đối chiếu dữ liệu công bố với bản ghi trên blockchain để phát hiện nếu có sự can thiệp sau khi bầu cử kết thúc.
\end{itemize}
